\documentclass[11pt,a4paper]{article}

\usepackage[utf8]{inputenc}
\usepackage[T1]{fontenc}
\usepackage{lmodern}
\usepackage[margin=22mm,headheight=15pt]{geometry}
\usepackage{microtype}
\usepackage{parskip}
\usepackage{enumitem}
\usepackage{booktabs}
\usepackage{longtable}
\usepackage{tabularx}
\usepackage{xcolor}
\usepackage{hyperref}
\usepackage{fancyhdr}
\usepackage{listings}
\usepackage{amsmath}
\usepackage{amssymb}

\definecolor{AionBlue}{HTML}{126E82}
\definecolor{AionNavy}{HTML}{17324D}
\definecolor{AionLight}{HTML}{EAF5F7}
\definecolor{AionGrey}{HTML}{5D6872}
\definecolor{CodeBg}{HTML}{F4F6F7}
\definecolor{WarningBg}{HTML}{FFF4DD}
\definecolor{RiskRed}{HTML}{A63A3A}

\hypersetup{
  colorlinks=true,
  linkcolor=AionNavy,
  urlcolor=AionBlue,
  pdftitle={AION Native English Intelligence Architecture and Operations},
  pdfauthor={Tessaris / AION Engineering}
}

\pagestyle{fancy}
\fancyhf{}
\lhead{AION Native English Intelligence}
\rhead{Architecture and Operations}
\cfoot{\thepage}

\setlist[itemize]{leftmargin=*,itemsep=2pt,topsep=3pt}
\setlist[enumerate]{leftmargin=*,itemsep=2pt,topsep=3pt}
\renewcommand{\arraystretch}{1.18}

\lstdefinestyle{aioncode}{
  backgroundcolor=\color{CodeBg},
  basicstyle=\ttfamily\small,
  breaklines=true,
  frame=single,
  rulecolor=\color{AionLight},
  columns=fullflexible,
  keepspaces=true,
  showstringspaces=false
}
\lstset{style=aioncode}

\newcommand{\statusbox}[1]{%
  \begin{center}
  \fcolorbox{AionBlue}{AionLight}{\parbox{0.92\textwidth}{#1}}
  \end{center}}

\newcommand{\warningbox}[1]{%
  \begin{center}
  \fcolorbox{RiskRed}{WarningBg}{\parbox{0.92\textwidth}{#1}}
  \end{center}}

\title{\textcolor{AionNavy}{\Huge AION Native English Intelligence}\\[4mm]
\Large Architecture, Learning Governance, Storage and Operations}
\author{Tessaris / AION Engineering}
\date{Architecture baseline: 7 October 2026}

\begin{document}
\maketitle

\begin{abstract}
This document records the design and operational state of AION's restored native English
intelligence layer. The system combines a verified local lexical and semantic store, a
model-free student and examiner, a governed local Qwen teacher, a fifteen-layer curriculum,
and a native symbolic conversation pipeline that parses, retrieves, reasons, plans and
realises English without next-token prediction. It also
documents the removable-media deployment, database schema, source files, service controls,
tests, recovery procedures, known limitations, and the conditions under which AION may claim
progress. The central design rule is that Qwen may teach curriculum material but may not
answer AION's examinations, compose or polish conversational answers, participate in answer
checking, or award AION competence.
\end{abstract}

\statusbox{\textbf{Operational snapshot, 7 October 2026.} The SD-resident native store
contains 147,306 verified lexemes, 982,850 semantic relations, 223 separated teacher lessons
and 2,438 closed-book attempts at 99.1\% recorded accuracy. AION has completed the original
seven-layer foundation and is on layer 9 of 15, \texttt{paragraph\_comprehension}, with one of
two required qualifying passes. The native conversation ledger contains 241 turns and six
self-improvement cycles. The service is healthy, the SD store is mounted with 13.78 GiB free,
and native conversational responses declare no Qwen, model, cloud, external polisher or
next-token-prediction route. These are dated measurements, not permanent capability claims.}

\tableofcontents
\newpage

\section{Purpose and Scope}

The goal is to make English knowledge part of AION's own durable cognitive substrate rather
than making every English question a disguised call to a generative model. The implemented
system restores the earlier lexical-conceptual intention---definitions, relationships,
context and continuous self-training---while connecting it to the current Learning Centre,
governance gates, desktop interface and local-model infrastructure.

This implementation is intentionally narrower than a claim of human-level language
understanding. It establishes and tests lexical grounding, semantic relations, contrasts,
context completion, mixed transfer, sentence semantics, discourse reasoning and a sequence
of advanced free-response layers. It also supplies a usable native conversation surface.
The remaining limitation is breadth: the symbolic parser, verified fact base and native
grammar must continue expanding before unrestricted general conversation can be claimed.

\subsection{What ``native'' means here}

In this document, \emph{native English intelligence} means that AION can retrieve, compare,
associate and answer supported English questions using its own local structured memory and
deterministic reasoning path. For a native examination:

\begin{itemize}
  \item no cloud service is used;
  \item Qwen is not called by the student path;
  \item the expected answer is withheld until after the student response;
  \item grading is deterministic;
  \item reinforcement is subject to the Cognitive Authority Unit (CAU); and
  \item the evidence is bounded to the tested subskills and task family.
\end{itemize}

\subsection{What the implementation does not claim}

The system does not yet prove unrestricted conversation, broad commonsense or human-level
pragmatic judgement. It now performs controlled compositional parsing, reference resolution,
paragraph inference and pragmatic tests, but success is evidence only for the named sealed
families. The architecture preserves this distinction in promotion logic, provenance and
explicit abstention.

\section{System Architecture}

The architecture is a governed teacher--student system with four independent roles.

\begin{center}
\begin{tabularx}{\textwidth}{>{\bfseries}p{0.18\textwidth} X}
\toprule
Role & Responsibility \\
\midrule
Qwen teacher & Receives the goal, complete curriculum, current stage, verified study targets,
native-memory statistics and previous closed-book results. It creates explanations,
contrasts, examples and guided study material. \\
AION student & Answers supported exercises from the native SQLite lexicon and semantic
relations. It never receives the expected-answer field. \\
Native examiner & Generates deterministic tasks, withholds the answer key, compares the
student response with verified ground truth, and records the result. \\
Learning Centre & Decides which subject contracts run, accepts only bounded evidence,
tracks covered subskills and controls broader curriculum advancement. \\
\bottomrule
\end{tabularx}
\end{center}

The runtime flow is:

\begin{enumerate}
  \item Load the current curriculum state and native lexical statistics.
  \item Select verified terms appropriate to the active curriculum stage.
  \item Present Qwen with the global goal, full curriculum, stage, targets and prior results.
  \item Store Qwen's response as a \emph{teacher lesson}; it is explicitly marked as not being
        examination evidence.
  \item Close the teacher interaction.
  \item Generate a new native examination with a withheld answer key.
  \item Let AION answer from its local lexicon and relation graph.
  \item Grade deterministically and consult CAU before reinforcement.
  \item Update stage history. Two qualifying stage passes are required before promotion.
  \item Expose the bounded result to the Learning Centre and service status.
\end{enumerate}

\subsection{Why this resembles distillation}

The design is curriculum-driven capability transfer. It resembles model distillation because
a more capable teacher shapes a bounded student. It differs from conventional neural
distillation because knowledge is currently transferred into structured memory, semantic
relations, practice history and verified skills rather than directly into smaller neural
weights. A future fine-tuning phase could consume verified records, but that is not part of
the present implementation.

\section{Core Native Lexical Runtime}

The principal runtime is:

\begin{quote}
\path{backend/modules/aion_cognition/native_lexical_intelligence.py}
\end{quote}

It defines \texttt{NativeLexicalIntelligence}, schema version 3. The class owns storage,
ingestion, lookup, semantic relation queries, native question answering, task generation,
student answering, grading, governed reinforcement and statistics.

\subsection{Supported native operations}

\begin{longtable}{p{0.25\textwidth} p{0.67\textwidth}}
\toprule
Operation & Behaviour \\
\midrule
\endhead
\texttt{lookup} & Retrieves the normalized lemma, part of speech, definitions, examples,
source, verification state and strongest semantic relations. \\
\texttt{related} & Retrieves synonym, antonym, hypernym, hyponym or derivational relations,
ordered by absolute relation strength. \\
\texttt{compare} & Compares two known entries using direct relations and definition-token
overlap. \\
\texttt{context} & Ranks a word's known definitions against the surrounding sentence tokens. \\
\texttt{answer} & Handles bounded requests including definitions, synonyms, antonyms,
comparisons and the explicit context query form. \\
\texttt{generate\_tasks} & Creates definition, synonym, antonym, context-completion or mixed
tasks from verified entries. \\
\texttt{answer\_task} & Executes the student path using only the prompt, choices and permitted
metadata. The expected answer is absent. \\
\texttt{run\_practice\_cycle} & Generates tasks, obtains student answers, grades them, records
attempts and reinforces only correct CAU-authorized results. \\
\texttt{record\_teacher\_lesson} & Stores teacher material in a separate table and labels it as
not counting toward exam evidence. \\
\bottomrule
\end{longtable}

\subsection{Normalization and scoring}

Lemmas are lower-cased, underscores are converted to spaces, and repeated whitespace is
collapsed. Definition matching and context completion use token overlap against stored
definitions and examples. Synonym and antonym questions use verified graph edges. This is a
transparent baseline rather than a neural semantic encoder; its strengths and weaknesses can
therefore be inspected exactly.

\section{Database Design}

The active database is:

\begin{quote}
\path{/Volumes/TESSARISQ38/TessarisLearning/AION/native_lexical_intelligence.sqlite3}
\end{quote}

SQLite is used because it is local, portable, transactional, inspectable and does not require
a database server. The SD-card runtime deliberately uses rollback journal mode
\texttt{DELETE} with \texttt{synchronous=FULL}; WAL is avoided because removable filesystems
may not provide reliable shared-memory sidecars. Internal-disk fallback uses WAL with normal
synchronisation. The database now contains both lexical-learning and native-conversation
structures.

\subsection{Lexemes}

The \texttt{lexemes} table has one row per normalized lemma. Important columns are:

\begin{itemize}
  \item \texttt{lemma}: primary key;
  \item \texttt{part\_of\_speech};
  \item \texttt{definition} and \texttt{definitions\_json};
  \item \texttt{examples\_json};
  \item \texttt{source} and \texttt{source\_ref};
  \item \texttt{verified};
  \item \texttt{reinforcement\_count}; and
  \item creation and update epochs.
\end{itemize}

\subsection{Semantic relations}

The \texttt{relations} table is keyed by source lemma, relation type and target lemma. Supported
relation types currently include synonym, antonym, hypernym, hyponym and derivation. Each edge
stores its strength, source and verification flag. An index supports efficient source-and-type
queries.

\subsection{Practice attempts}

The \texttt{practice\_attempts} table is the auditable closed-book record. It stores task ID,
family, lemma, prompt, student answer, expected answer, correctness, CAU authorization,
project family and timestamp. The expected answer is stored only after the student answer is
produced; it is not present in the student's input mapping.

\subsection{Teacher lessons}

The \texttt{teacher\_lessons} table deliberately separates teaching from evidence. It stores
the stage, lesson JSON, referenced lemmas, teacher model, count of verified references,
authority status and timestamp. A lesson can improve study exposure, but the row cannot by
itself raise a competency level.

\subsection{Metadata}

The \texttt{metadata} table records the schema version and last-ingest epoch. It provides a
small migration and operational anchor without mixing system metadata into lexical entries.

\section{Corpus and Provenance}

The restored base corpus is Princeton WordNet 3.0, imported through NLTK. The ingestion helper
\texttt{wordnet\_records} aggregates synsets into lemma records and derives definitions,
examples, synonyms, antonyms, hypernyms, hyponyms and derivational links.

The import snapshot produced:

\begin{center}
\begin{tabular}{lr}
\toprule
Verified lexemes & 147,306 \\
Semantic relations & 982,850 \\
Skipped records & 0 \\
\bottomrule
\end{tabular}
\end{center}

The downloaded source corpus is retained at:

\begin{quote}
\path{/Volumes/TESSARISQ38/TessarisLearning/AION/nltk_data}
\end{quote}

The generated SQLite database is the runtime dependency. NLTK is needed for rebuilding or
extending the WordNet import, but ordinary lookup and examination use the Python standard
library and SQLite.

\warningbox{WordNet provides broad lexical coverage, not complete modern conversational
English and not verified world knowledge. Some lemmas are obscure, historical, multiword or
domain-specific. Future portfolio sampling should include frequency bands and curated modern
usage sets so perfect performance is not dominated by internally familiar dictionary edges.}

\section{The Qwen Teacher}

The teacher orchestration is implemented in:

\begin{quote}
\path{backend/modules/aion_cognition/native_english_teacher.py}
\end{quote}

The teacher prompt contains five essential parts:

\begin{enumerate}
  \item the permanent goal;
  \item the complete ordered curriculum;
  \item the active stage;
  \item AION's current memory and prior closed-book results; and
  \item a bounded set of verified study targets and definitions.
\end{enumerate}

Qwen is asked to return a compact JSON lesson containing an objective, explanation, contrasts,
examples, guided practice, referenced lemmas and next focus. Teacher output is treated as a
proposal. Referenced terms are checked against the native lexicon, and the lesson is stored
separately from verified facts and examination evidence.

\subsection{Local model configuration}

The teacher uses the approved bounded local package:

\begin{itemize}
  \item model ID: \texttt{qwen3-8b-fast-local-q4};
  \item runtime adapter: \texttt{openai\_compatible\_bounded\_chat.v1};
  \item endpoint: \url{http://127.0.0.1:8897};
  \item context size: 4096;
  \item parallel slots: 1;
  \item maximum output tokens: 320;
  \item network fallback: disabled;
  \item tool execution: disabled; and
  \item finance arithmetic: forced through separately verified Tessaris tools.
\end{itemize}

The exact packaged weight is located under:

\begin{quote}
\path{/Volumes/TESSARISQ38/TessarisModelPackages/qwen3-8b-fast-local/weights/}
\end{quote}

The runtime normally requires a fully selectable verified installation. One carefully bounded
exception is implemented for an already-running server: if the SD warehouse is verified, the
model path exists, the only preflight failure is insufficient internal free-space reserve, and
the declared loopback health endpoint is responsive, AION may reuse that already-loaded exact
runtime. It may not use this exception for hash, warehouse, model-source, architecture or
memory failures.

\section{Curriculum and Promotion}

The staged curriculum is defined in \texttt{native\_english\_teacher.py}.

\begin{longtable}{p{0.05\textwidth} p{0.25\textwidth} p{0.36\textwidth} p{0.10\textwidth} p{0.10\textwidth}}
\toprule
No. & Stage ID & Focus & Minimum & Threshold \\
\midrule
\endhead
1 & \texttt{lexical\_foundations} & Word meaning and definition matching & 24 & 0.82 \\
2 & \texttt{semantic\_relations} & Synonyms and concept relations & 32 & 0.82 \\
3 & \texttt{semantic\_contrasts} & Antonyms and fine-grained contrasts & 32 & 0.82 \\
4 & \texttt{contextual\_meaning} & Meaning in unfamiliar sentence context & 40 & 0.84 \\
5 & \texttt{mixed\_transfer} & Mixed lexical transfer without teacher access & 48 & 0.86 \\
6 & \texttt{sentence\_semantics} & Entities, actions, properties, tense and negation & 64 & 0.88 \\
7 & \texttt{discourse\_reasoning} & Reference, ambiguity, argument and pragmatics & 80 & 0.90 \\
8 & \texttt{compositional\_syntax} & Complex clauses and compositional syntax & 80 & 0.90 \\
9 & \texttt{paragraph\_comprehension} & Paragraph meaning and causal sequence & 96 & 0.90 \\
10 & \texttt{unseen\_document\_qa} & Unseen-document reading and exact retrieval & 96 & 0.92 \\
11 & \texttt{dialogue\_tracking} & Multi-turn reference and instruction tracking & 96 & 0.92 \\
12 & \texttt{ambiguity\_pragmatics} & Implication and calibrated uncertainty & 96 & 0.90 \\
13 & \texttt{explanation\_generation} & Grounded explanation in AION's own words & 96 & 0.90 \\
14 & \texttt{error\_correction} & Grammar diagnosis and precise correction & 96 & 0.94 \\
15 & \texttt{cross\_domain\_transfer} & Cross-domain reasoning and reliability & 128 & 0.92 \\
\bottomrule
\end{longtable}

For a stage with threshold $\tau$ and minimum attempt count $N_{\min}$, a stage pass is:

\[
\mathrm{StagePass} =
  (N \geq N_{\min}) \land
  (U \geq N_{\min}) \land
  \left(\frac{C}{N} \geq \tau\right) \land
  E_{\mathrm{current}} \land W \land M_{0} \land A,
\]

where $U$ is the number of unique tasks, $E_{\mathrm{current}}$ means the current examination
also meets the threshold, $W$ means the answer key was withheld, $M_0$ means no model or cloud
route was used by the student, and $A$ is CAU authorization. Layers 6--15 additionally require
the minimum number of source-disjoint attempts. Layers 8--15 require free responses rather
than answer selection. Promotion requires two consecutive qualifying passes. The teacher does
not participate in this calculation.

The version-3 migration preserves the completed seven-layer foundation. A prior state with
\texttt{curriculum\_complete=true} at layer 7 becomes
\texttt{foundation\_curriculum\_complete=true}, moves to layer 8 and reopens the expanded
curriculum without deleting previous evidence.

\section{Closed-Book Examination Guarantees}

Every generated \texttt{LexicalTask} contains an expected answer internally, but
\texttt{student\_view()} removes it. The student receives only the task ID, task family,
prompt, choices and---only when needed for relation queries---the lemma. Grading occurs after
\texttt{answer\_task()} returns.

The teacher cycle emits the following explicit invariants:

\begin{itemize}
  \item \texttt{teacher\_saw\_exam\_answer\_key = false};
  \item \texttt{teacher\_graded\_exam = false};
  \item \texttt{model\_used\_for\_student\_answers = false}; and
  \item \texttt{lesson\_counts\_as\_competence\_evidence = false}.
\end{itemize}

These flags are useful audit assertions, but the stronger guarantee is structural: the call to
the teacher finishes before the independent practice-cycle call begins, and the practice call
does not accept a teacher callback or lesson payload.

\section{Cognitive Authority and Reinforcement}

Native learning consults:

\begin{quote}
\path{backend/modules/aion_cognition/cau_authority.py}
\end{quote}

CAU can permit reading while denying reinforcement. This matters because AION should remain
able to answer from existing verified memory during an unstable state, while being prevented
from writing new reinforcement when coherence or stability conditions fail. The practice
record preserves whether learning authority was available for each attempt.

The current live snapshot reported high stability and allowed native lexical learning. Those
values are runtime conditions, not constants; denial must remain fail-closed if CAU is
unavailable.

\section{Learning Centre Integration}

The closed-book project authority is implemented in:

\begin{quote}
\path{backend/modules/hexcore/native_english_competency_portfolio.py}
\end{quote}

It provides eight versioned project families spanning definitions, synonyms, antonyms,
context and mixed transfer. A project is accepted only when it has at least 12 trials, a score
of at least 0.75, a withheld answer key, no student model use and CAU authorization. It emits
bounded evidence authorities:

\begin{itemize}
  \item \texttt{native\_lexical\_memory};
  \item \texttt{withheld\_answer\_key}; and
  \item \texttt{cau\_learning\_gate}.
\end{itemize}

The English project route in
\path{backend/modules/hexcore/progressive_competency_executor.py} now uses this portfolio
instead of the old Python-snippet English simulation. The executor also records the actual
trial count and authority returned by this portfolio rather than incorrectly labelling it as
a Python subprocess and security scan.

The English subject registry in
\path{backend/modules/hexcore/constitutional_north_star_mastery_registry.py} now includes:

\begin{itemize}
  \item lexical grounding;
  \item semantic relations;
  \item morphology;
  \item syntax;
  \item sentence semantics;
  \item reference resolution; and
  \item native question answering;
\end{itemize}

alongside the existing instruction, document, argument, ambiguity, correction, generation,
pragmatic and dialogue subskills. The migration in
\path{backend/modules/hexcore/progressive_competency_system.py} merges newly declared
subskills into persisted subject state without erasing existing evidence.

\section{Continuous Service Operation}

The persistent service is:

\begin{quote}
\path{backend/AION/system/aion_mastery_curriculum_service.py}
\end{quote}

It starts a single-worker background teacher pool before the longer Learning Centre startup
work. This prevents a slow Qwen response from blocking the main curriculum loop. The default
teacher interval is 300 seconds, with a hard minimum of 60 seconds. Results appear in the
service status under \texttt{native\_english\_teacher}.

The service is supervised by:

\begin{quote}
\path{backend/AION/system/aion_heartbeat.py}
\end{quote}

and on the current Mac runs under the launch agent
\texttt{com.tessaris.aion.heartbeat}. The supervisor restarts the mastery service if its
process terminates.

\subsection{Environment controls}

\begin{longtable}{p{0.36\textwidth} p{0.54\textwidth}}
\toprule
Variable & Meaning \\
\midrule
\endhead
\texttt{AION\_NATIVE\_ENGLISH\_TEACHER} & Set to \texttt{0}, \texttt{false} or
\texttt{False} to disable continuous teacher dispatch. Default is enabled. \\
\texttt{AION\_NATIVE\_ENGLISH\_TEACHER\_INTERVAL} & Seconds between dispatches. Default 300;
minimum 60. \\
\texttt{AION\_NATIVE\_LEXICON\_DB} & Explicit absolute database path. This takes precedence
over automatic SD discovery. \\
\texttt{AION\_REPO\_ROOT} & Repository root used by persistent services. \\
\texttt{TESSARIS\_LOCAL\_MODEL\_ROOTS} & Optional colon-separated additional model-vault
scan roots. \\
\bottomrule
\end{longtable}

\section{Question and Dashboard Integration}

The dashboard action layer is:

\begin{quote}
\path{backend/modules/aion_dashboard/actions.py}
\end{quote}

Its \texttt{ask} action now calls native lexical answering rather than the prior canned or
random query path. Its \texttt{teach} action runs native closed-book practice instead of
passing the expected answer directly to a shallow evaluator.

The current desktop/local-node route is:

\begin{quote}
\path{backend/modules/local_node/local_node_runtime.py}
\end{quote}

The older general desktop chat route gives native lexical lookup first refusal and may send an
unhandled question to a selected local model. That route must not be confused with the new
native-English conversation contract. The dedicated
\texttt{/english-intelligence/conversation/ask} route never falls back to a model: it answers,
asks for clarification or abstains. Native AION speech and model-assisted desktop chat must
remain visibly and technically distinct.

As a verified example, the prompt ``What does inference mean?'' was answered directly from the
WordNet-backed native entry and returned its definition and semantic relations without calling
Qwen.

\section{Native English Conversation Runtime}

The new model-free speaking path is implemented in:

\begin{quote}
\path{backend/modules/aion_cognition/native_english_conversation.py}
\end{quote}

The class \texttt{NativeEnglishConversation} is deliberately not an LLM wrapper. It has no
prompt builder, token sampler, completion client, embedding API, cloud client or polishing
callback. A question passes through six explicit stages:

\begin{enumerate}
  \item \textbf{Parse}: normalize text and identify question intent, slots, tokens, tense,
        negation, candidate entities, action, subject terms and object terms.
  \item \textbf{Resolve context}: when a slot contains a reference such as ``it'', ``this''
        or ``they'', inspect the latest answered turn in the same session and bind the
        reference to a prior term, subject or object. Unresolved references cause
        clarification rather than guessing.
  \item \textbf{Retrieve}: consult verified lexical entries, verified semantic relations or
        verified structured facts. Unverified teacher prose is not an answer source.
  \item \textbf{Plan}: choose a symbolic answer act and assemble a proposition from retrieved
        meanings and evidence identifiers.
  \item \textbf{Realise}: convert the proposition into a grammatical surface sentence with
        explicit rules for capitalization, punctuation, relation lists, comparisons,
        explanations and factual clauses.
  \item \textbf{Check}: require evidence, basic grammatical surface validity, balanced quoted
        text and the absence of any model/cloud route. Unsupported questions are declined.
\end{enumerate}

This is meaning-plan-to-language generation. It does not predict a next token from a probability
distribution. Selection still occurs---AION selects an intent, evidence set, answer plan and
grammar rule---but those decisions are inspectable symbolic operations.

\subsection{Native response contract}

Every answer returns the parsed frame, status, confidence, evidence, checks and a generation
receipt. The receipt is designed to be machine-auditable:

\begin{lstlisting}
{
  "method": "symbolic_meaning_plan_and_rule_grammar",
  "aion_generated": true,
  "next_token_prediction": false,
  "qwen_used": false,
  "model_used": false,
  "cloud_used": false,
  "external_polisher_used": false
}
\end{lstlisting}

These flags are backed by structure: the native conversation module imports the lexical
runtime and SQLite only. It does not import the local-model runtime. A future integration must
not weaken this boundary silently; tests should fail if a model route is introduced.

\subsection{Supported conversational functions}

\begin{longtable}{p{0.27\textwidth} p{0.64\textwidth}}
\toprule
Function & Current behaviour \\
\midrule
\endhead
Definition & Produces ``The word X means Y'' from a verified lexical entry. \\
Synonym/antonym & Retrieves verified relation edges and realises a bounded English list. \\
Comparison & Retrieves both definitions, direct relations and shared semantic material before
constructing the comparison. \\
Known-about query & Summarizes a verified lexical concept and up to four relations. \\
Structured fact query & Answers supported who/what/where/when/why forms from verified fact
records and qualifiers. \\
Session reference & Resolves selected pronouns and demonstratives from the latest answered turn
within the same conversation session. \\
Clarification & Requests a subject and requested information when the parse is ambiguous. \\
Abstention & Says that verified knowledge is insufficient when the question is understood but
the evidence store cannot support an answer. \\
Audit & Persists question, parse, answer, status, confidence, evidence and checks. \\
\bottomrule
\end{longtable}

The live two-turn acceptance check was:

\begin{quote}
\textbf{Question 1:} Define calm.\\
\textbf{Native answer:} The word `calm' means (of weather) free from storm or wind.\\[2mm]
\textbf{Question 2:} What does it mean?\\
\textbf{Native answer:} The word `calm' means (of weather) free from storm or wind.
\end{quote}

The second result records \texttt{context.used=true}, the source turn identifier and
\texttt{resolved\_to=calm}. Both answers report no model, Qwen, cloud, polisher or token
prediction.

\subsection{Verified fact interface}

\texttt{remember\_fact} stores a subject, predicate, object, optional place/time/cause
qualifiers, provenance and verification state. Mutation is CAU-gated. The public HTTP API does
not expose an unrestricted ``mark this fact verified'' endpoint; this avoids allowing an
arbitrary caller to poison AION's native truth store. Trusted subsystem adapters must supply
source identity and verification before calling the internal method.

At the dated snapshot the structured fact count is zero, so broad factual questions correctly
abstain. This is a capability boundary, not a defect to hide: useful general answers require
verified business, world, project and experiential facts to be bridged into this table or into
another governed retriever implementing the same evidence contract.

\section{Native Self-Improvement Programme}

\texttt{self\_improvement\_cycle} provides continuous improvement without using Qwen or any
other model as answerer, grader or editor. It has two initial mechanisms:

\begin{enumerate}
  \item retry open failures after memory changes and resolve them only when a new native answer
        has verified evidence; and
  \item generate paraphrases of verified definition questions and require both forms to return
        the same answered state and evidence identity.
\end{enumerate}

The cycle is CAU-gated and records its own receipt. It returns
\texttt{model\_used=false}, \texttt{cloud\_used=false} and
\texttt{next\_token\_prediction=false}. Invalid internally generated probes are discarded
instead of becoming permanent curriculum failures. The cycle is invoked after each teacher
cycle and is also available as an explicit local API operation.

\warningbox{``Self-improvement'' must not mean self-certification. AION may discover a gap,
retry a question, test paraphrase consistency and reinforce an independently verified result.
It may not convert its own unsupported sentence into truth merely because it generated the
sentence repeatedly. This prevents autonomous error amplification.}

The current design should expand through versioned mechanisms:

\begin{itemize}
  \item cluster failures by missing intent, grammar rule, reference type or knowledge source;
  \item create remedial exercises from verified facts without exposing their answers;
  \item test adversarial paraphrases, negation and minimally changed sentences;
  \item schedule delayed retention at 1, 7, 30 and 90 days;
  \item promote a new parser or grammar rule only after sealed regression tests pass; and
  \item roll back a rule version when it reduces performance on unrelated held-out families.
\end{itemize}

\section{Advanced Examination Engine}

Advanced model-free examinations are implemented in:

\begin{quote}
\path{backend/modules/aion_cognition/native_english_reasoning.py}
\end{quote}

The engine contains deterministic task construction and a student solver that receives only
the prompt, response mode and non-answer metadata. Expected answers remain inside the examiner
until the student response is returned. Sentence-semantics and discourse families use selected
responses where appropriate. Layers 8--15 use free response and mark every task as
source-disjoint from the WordNet-derived lexical curriculum.

Eight advanced families currently provide 128 unique sealed situations each, for 1,024
advanced free-response situations:

\begin{itemize}
  \item complex clauses and modifier scope;
  \item paragraph cause and sequence;
  \item unseen note reading with constrained calculation;
  \item revised instructions across dialogue turns;
  \item implication, ambiguity and explicit uncertainty;
  \item evidence-grounded explanation;
  \item tense and agreement correction; and
  \item reliability reasoning across science, history, engineering, ecology, medicine,
        finance, law and computing contexts.
\end{itemize}

Task identifiers hash the family, prompt, response space and expected answer. Mixed-radix
selection and case variation avoid the short modulo cycles that would otherwise create many
nominal attempts from a small number of semantic situations. Promotion counts distinct task
identifiers, not just raw attempts.

\section{Conversation Storage Schema}

The native conversation tables share the SD-resident SQLite store so language memory remains
portable with the lexical substrate:

\begin{longtable}{p{0.34\textwidth} p{0.56\textwidth}}
\toprule
Table & Purpose \\
\midrule
\endhead
\texttt{native\_english\_facts} & Verified subject--predicate--object propositions,
qualifiers, provenance and timestamps. \\
\texttt{native\_english\_turns} & Session, question, parse JSON, answer, status, confidence,
evidence, checks and timestamp for every native turn. \\
\texttt{native\_english\_failures} & Stable failure identity, question, reason, parse, retry
count and lifecycle status. \\
\texttt{native\_english\_improvement\_cycles} & Examined/resolved failure counts,
paraphrase-consistency results, authority state and cycle details. \\
\bottomrule
\end{longtable}

Conversation writes use short transactions and the same 30-second busy timeout as the lexical
runtime. This is important because the SD card may be accessed simultaneously by the Learning
Centre service and the desktop API. Rollback journalling avoids the previous WAL/SHM failure
mode on removable media.

\section{English Intelligence API and Panel}

The HTTP surface and visual progress panel are implemented in:

\begin{quote}
\path{backend/api/aion_english_intelligence_router.py}
\end{quote}

\begin{longtable}{p{0.46\textwidth} p{0.44\textwidth}}
\toprule
Route & Function \\
\midrule
\endhead
\texttt{GET /api/aion/english-intelligence/progress} & Read-only health, curriculum, storage,
exam, Learning Centre and conversation metrics. \\
\texttt{GET /api/aion/english-intelligence/panel} & Dark-theme operational panel with live
refresh, curriculum layers and native conversation tester. \\
\texttt{POST /api/aion/english-intelligence/conversation/ask} & Run one native conversation
turn with a question and session identifier. \\
\texttt{POST /api/aion/english-intelligence/conversation/improve} & Run one governed native
self-improvement cycle. \\
\bottomrule
\end{longtable}

The panel explicitly displays the proof boundary beneath every answer: confidence, verified
evidence count and the statement ``Native Aion---no Qwen---no cloud---no token prediction.''
It is intended as both a user test surface and an operational diagnostic.

\section{Restored Legacy Components}

The earlier Phase 41/advanced-cognition components remain for compatibility, but key defects
were repaired:

\begin{itemize}
  \item \path{backend/modules/aion_language/query_resonance_api.py} now reads the same
        \path{data/lexicons/language_resonance_matrix.json} location written by the exporter
        and unwraps the export's \texttt{matrix} object.
  \item \path{backend/modules/aion_language/thesaurus_linker.py} now uses deterministic
        synonym and antonym strengths rather than time-varying sine values.
  \item \path{backend/modules/aion_cognition/advanced_cognition/lexical_task_runner.py} now
        routes compatibility task sets through the governed native examiner. It no longer
        picks synonym answers randomly, compares a definition directly with its lemma string,
        or manufactures random SQI and stability values.
\end{itemize}

\section{SD-Card Deployment}

\subsection{What is stored on the SD card}

The SD card currently contains two different classes of assets:

\begin{enumerate}
  \item \textbf{Teacher model assets}: the verified Qwen package and its warehouse manifest
        under \path{/Volumes/TESSARISQ38/TessarisModelPackages/}.
  \item \textbf{Native English assets}: the SQLite intelligence store and retained NLTK corpus
        under \path{/Volumes/TESSARISQ38/TessarisLearning/AION/}.
\end{enumerate}

The AION codebase, curriculum definitions, most competency records, teacher-cycle state and
most other AION memories remain in the Mac repository. The SD card is therefore an attached
language/model volume, not the whole AION system.

\subsection{Automatic database resolution}

Database resolution follows this order:

\begin{enumerate}
  \item use \texttt{AION\_NATIVE\_LEXICON\_DB} when explicitly configured;
  \item scan mounted volumes for
        \path{TessarisLearning/AION/native_lexical_intelligence.sqlite3};
  \item otherwise use \path{data/aion/native_lexical_intelligence.sqlite3} under the repository.
\end{enumerate}

The SD deployment was chosen when internal temporary space was critically low. At the dated
snapshot used for this document, the generated database was approximately 192 MB, the retained
NLTK corpus approximately 36 MB, the SD volume had approximately 14 GiB available, and macOS
later reported approximately 7.6 GiB internally available after system cleanup. Capacity is
dynamic and must be checked before copying the database back to the Mac.

\warningbox{Current operational risk: if the SD card is absent and no explicit database path
is set, resolution can fall back to a new or existing internal repository database. That can
look like an empty AION rather than a disconnected memory volume. A fail-loud mount manifest
and user-visible ``English memory unavailable'' status should be implemented before production
deployment.}

\subsection{Mount and health checks}

The following read-only checks establish whether the expected assets are visible:

\begin{lstlisting}[language=bash]
test -f /Volumes/TESSARISQ38/TessarisLearning/AION/native_lexical_intelligence.sqlite3
test -f /Volumes/TESSARISQ38/TessarisModelPackages/qwen3-8b-fast-local/weights/Qwen_Qwen3-8B-Q4_K_M.gguf
curl --fail --silent http://127.0.0.1:8897/health
df -h / /Volumes/TESSARISQ38
\end{lstlisting}

The model health endpoint proves only that a server responds. Runtime code additionally checks
the selected model profile, loopback restriction, context and slot configuration, package path
and exact model hash.

\section{Backup, Portability and Recovery}

\subsection{Backup policy}

The database is portable but should not exist as a single unbacked-up intelligence volume.
At minimum, retain:

\begin{itemize}
  \item a SQLite-consistent backup of the database;
  \item a SHA-256 checksum for the backup;
  \item the teacher state and cycle ledger from
        \path{data/aion/native_english_teacher/};
  \item the relevant progressive competency state and evidence artifacts; and
  \item the model package and warehouse manifest independently of the lexical database.
\end{itemize}

Use SQLite's backup operation while services are live rather than copying only the main file
while write-ahead-log files may contain committed pages:

\begin{lstlisting}[language=bash]
sqlite3 /Volumes/TESSARISQ38/TessarisLearning/AION/native_lexical_intelligence.sqlite3 \
  ".backup '/path/to/backup/native_lexical_intelligence.snapshot.sqlite3'"
shasum -a 256 /path/to/backup/native_lexical_intelligence.snapshot.sqlite3
\end{lstlisting}

The placeholder destination must be replaced with a real backup volume; it must not point to
the same physical SD card.

\subsection{Moving the intelligence store to a Mac or robot}

There is no requirement to retrain from zero. A safe migration is:

\begin{enumerate}
  \item pause the native teacher service;
  \item create and verify a SQLite backup;
  \item copy the backup to the new device's durable data location;
  \item set \texttt{AION\_NATIVE\_LEXICON\_DB} to the new absolute path;
  \item start AION and compare lexeme, relation, lesson and attempt counts;
  \item run a read-only definition query;
  \item run a small closed-book test; and
  \item resume continuous teaching only after CAU allows learning.
\end{enumerate}

The code and the data should be versioned separately. Repository checkout alone does not carry
the live 192 MB database, and copying the database alone does not carry the Learning Centre
logic.

\subsection{Recovery checks}

After restoration, verify at least:

\begin{itemize}
  \item SQLite opens and reports schema version 3;
  \item lexeme and relation counts match the backup receipt;
  \item practice-attempt counts and teacher-lesson counts are present;
  \item a known definition query succeeds without a model;
  \item Qwen remains local and network fallback remains disabled;
  \item a new exam reports \texttt{expected\_withheld = true}; and
  \item no stage promotion occurs solely from restored teacher lessons.
\end{itemize}

\section{Testing and Verification}

The focused automated tests are located in:

\begin{quote}
\path{backend/tests/test_native_lexical_intelligence.py}
\end{quote}

and:

\begin{quote}
\path{backend/tests/test_native_english_conversation.py}\\
\path{backend/tests/test_native_english_progress.py}\\
\path{backend/tests/test_progressive_competency_system.py}
\end{quote}

They verify that:

\begin{itemize}
  \item definitions, synonyms and comparisons work without a model;
  \item expected answers are missing from the student view;
  \item model and cloud usage flags remain false during practice;
  \item CAU denial blocks reinforcement but not reading;
  \item Qwen receives the goal and curriculum; and
  \item Qwen lessons do not count as examination evidence or grading authority;
  \item all eight advanced banks contain 128 unique free-response tasks;
  \item advanced answers are model-free, cloud-free and source-disjoint;
  \item the completed version-2 foundation migrates into version 3 without false completion;
  \item native answers carry evidence and the no-model generation receipt;
  \item unsupported factual questions produce abstention rather than fabrication;
  \item verified structured facts can be retrieved and realised as English;
  \item session references resolve from AION's own prior turn; and
  \item self-improvement uses verified paraphrase consistency without a model.
\end{itemize}

On 7 October 2026, 56 focused tests passed. Python compilation checks and
\texttt{git diff --check} also passed for the changed implementation files. Live acceptance
tests additionally exercised the HTTP conversation route, model-free receipt and two-turn
reference resolution against the SD-resident store.

\section{Operational Monitoring}

Important live artifacts are:

\begin{longtable}{p{0.42\textwidth} p{0.48\textwidth}}
\toprule
Path & Purpose \\
\midrule
\endhead
\path{data/aion/native_english_teacher/state.json} & Current stage, cycle count, consecutive
passes, last exam and last teacher error/model. \\
\path{data/aion/native_english_teacher/cycles.jsonl} & Append-only teacher-cycle and exam audit. \\
\path{results/aion_mastery_curriculum_service_status.json} & Persistent service status,
including the background native English teacher lane. \\
\path{backend/modules/hexcore/data/progressive_competency/state.json} & Curriculum subjects,
contracts, evidence and blockers. \\
\path{backend/logs/aion_heartbeat.log} & Service supervision and restart events. \\
\bottomrule
\end{longtable}

The most useful dashboard metrics are:

\begin{itemize}
  \item current stage and time since last successful cycle;
  \item teacher available/error and selected model;
  \item verified lexeme and relation counts;
  \item lessons stored, explicitly separated from exam evidence;
  \item closed-book attempts today, accuracy and task-family distribution;
  \item model/cloud usage flags for native examinations;
  \item CAU-authorized versus denied attempts;
  \item consecutive qualifying stage passes;
  \item blocked native exam families; and
  \item SD-card mount, free space, database health and backup age.
  \item native conversation turns, answer rate, open failures and improvement cycles;
  \item evidence count and generation receipt for each native answer; and
  \item unresolved Learning Centre contracts independently of the native curriculum state.
\end{itemize}

At the dated snapshot, the native English service itself has no blocker, while the broader
Learning Centre reports one open English executor-capability blocker on a lexical-grounding
knowledge-test contract. This distinction matters: a healthy native curriculum does not erase
an independently governed competency contract.

\section{Security, Integrity and Trust Boundaries}

\subsection{Teacher-output boundary}

Qwen-generated material is not automatically promoted to verified truth. It is stored as
study material, tagged with the teacher model and cross-referenced against known lemmas. The
system should continue moving toward claim-level verification for teacher examples and
explanations before allowing them to become new factual relations.

\subsection{Prompt boundary}

The teacher receives only bounded verified definitions, curriculum state and previous result
summaries. It is told that it is not the student, examiner or competence authority. The exam
is generated after the teacher call closes.

\subsection{Storage boundary}

The SD model warehouse is verified against its manifest and model SHA-256. The SQLite database
currently relies on filesystem access controls and SQLite consistency. Production hardening
should add signed backup receipts, a database checksum manifest, mount identity and optional
encryption at rest.

\subsection{Capability-claim boundary}

Learning Centre evidence is attached only to named subskills. One project family cannot be
reused indefinitely as novel evidence. Exhaustion requires a new portfolio version instead of
manufacturing progress from repeated identical work.

\section{Known Limitations}

\begin{enumerate}
  \item \textbf{Controlled capability is not unrestricted fluency.} The native parser and
        grammar cover declared intents and forms; unseen constructions may require
        clarification or abstention.
  \item \textbf{Verified world knowledge is sparse.} The lexical store is large, but the
        structured fact table is initially empty. Broad factual conversation requires
        governed bridges from AION's other memories and real outcomes.
  \item \textbf{Advanced tests are synthetic and sealed, not yet broad public benchmarks.}
        They prove the implemented operations and teacher separation, but should be augmented
        with licensed unseen documents and independent annotations.
  \item \textbf{Self-improvement cannot invent a truth signal.} It can diagnose, replay and
        verify consistency; new factual knowledge still needs evidence external to the answer
        being evaluated.
  \item \textbf{Reference resolution is intentionally conservative.} It currently uses the
        most recent answered turn and selected reference words rather than a complete discourse
        model.
  \item \textbf{Sampling quality still needs frequency control.} Lexical probes can over-sample
        obscure WordNet lemmas; modern frequency bands should be installed.
  \item \textbf{The SD card remains a single operational dependency.} Automatic independent
        backup, signed restore receipts and fail-loud volume identity remain required.
  \item \textbf{Qwen lessons remain separate study proposals.} They do not directly rewrite
        the native parser, grammar or verified fact store.
\end{enumerate}

\section{Required Next Engineering Layers}

\subsection{Layer 1: production storage safety}

Add an SD mount manifest, explicit configured database path, fail-loud health state, periodic
SQLite backup, checksum receipts, backup-age metric and restore drill. This should precede
substantial accumulation of unique learned material.

\subsection{Layer 2: governed knowledge bridges}

Connect project, business, scientific, experiential and AKG memories through adapters that
emit verified subject--predicate--object facts with provenance. Do not expose a generic public
endpoint that allows a caller to label arbitrary prose as verified.

\subsection{Layer 3: broader parser and grammar}

Expand modality, aspect, quantification, relative clauses, ellipsis, multi-reference chains,
register and longer response planning. Every new rule should have minimal-pair, negation and
regression portfolios before activation.

\subsection{Layer 4: delayed autonomous improvement}

Implement 1-, 7-, 30- and 90-day retention queues, failure clustering, versioned grammar-rule
promotion, regression detection and rollback. Improvement receipts must identify the external
truth or sealed authority that justified each change.

\subsection{Layer 5: external transfer}

Use unseen licensed or public text not used in corpus construction. Partition sources by
document and time so the test genuinely measures transfer. Maintain separate development,
validation, closed-book examination and delayed-retention sets.

\subsection{Layer 6: native response breadth}

Add multi-paragraph planning, safe summaries, explanations with multiple evidence items,
counterfactual comparison and calibrated alternative interpretations. These must remain
symbolic/native unless a separately labelled model-assisted mode is intentionally introduced;
model output must never be presented as native AION speech.

\section{Operator Runbook}

\subsection{Normal startup}

\begin{enumerate}
  \item Mount \texttt{TESSARISQ38}.
  \item Confirm the lexical database and Qwen weight exist.
  \item Confirm the Qwen loopback endpoint or allow the verified runtime manager to start it.
  \item Start or verify \texttt{com.tessaris.aion.heartbeat}.
  \item Inspect the mastery-service status and native teacher state.
  \item Confirm the last teacher error is null and the last native exam reports no model use.
\end{enumerate}

\subsection{If the SD card is disconnected}

Do not treat an empty internal database as loss of intelligence and do not re-import into a
second store without deciding which copy is authoritative. Pause the teacher, remount the card,
verify the database path and counts, then resume. If the card has failed, restore the latest
SQLite-consistent backup and its checksum receipt.

\subsection{If Qwen is unavailable}

Native lookup and closed-book examination can still operate from the lexical database. The
teacher cycle records an error but does not fabricate a lesson. Resolve model mount, package
verification, loopback endpoint and runtime-profile failures; do not enable network fallback
as a shortcut.

\subsection{If CAU denies learning}

Existing verified memory remains readable. New attempts can be recorded, but they must not
reinforce memory or qualify for stage progression. Investigate the reported CAU denial reason
and restore stable conditions rather than bypassing authority.

\subsection{If scores suddenly fall}

Check database path and counts first, then stage/task family, SD health, corpus integrity,
answer-key withholding, recent code changes and sampling distribution. Preserve failed
attempts: they are curriculum signals and should not be deleted merely to recover a high score.

\section{File and Folder Reference}

\begin{longtable}{p{0.48\textwidth} p{0.43\textwidth}}
\toprule
File or folder & Responsibility \\
\midrule
\endhead
\path{backend/modules/aion_cognition/native_lexical_intelligence.py} & Native memory,
relations, question answering and closed-book practice. \\
\path{backend/modules/aion_cognition/native_english_teacher.py} & Curriculum, Qwen teacher,
stage progression and teacher/exam separation. \\
\path{backend/modules/aion_cognition/native_english_reasoning.py} & Sealed sentence,
discourse and advanced free-response examination banks and native solvers. \\
\path{backend/modules/aion_cognition/native_english_conversation.py} & Symbolic question
parsing, context resolution, retrieval, answer planning, grammar realisation, checking,
conversation memory and model-free self-improvement. \\
\path{backend/modules/aion_cognition/native_english_progress.py} & Read-only health,
curriculum, storage, examination and Learning Centre projection. \\
\path{backend/api/aion_english_intelligence_router.py} & Progress API, visual panel, native
conversation route and explicit improvement route. \\
\path{backend/modules/hexcore/native_english_competency_portfolio.py} & Versioned Learning
Centre English project families. \\
\path{backend/modules/hexcore/progressive_competency_executor.py} & Routes English projects
and records bounded evidence. \\
\path{backend/modules/hexcore/progressive_competency_system.py} & Persistent subject catalogue,
contracts, evidence, blockers and subskill migration. \\
\path{backend/modules/hexcore/constitutional_north_star_mastery_registry.py} & Canonical English
subject and subskill declaration. \\
\path{backend/AION/system/aion_mastery_curriculum_service.py} & Continuous asynchronous teacher
dispatch and Learning Centre service status. \\
\path{backend/AION/system/aion_heartbeat.py} & Persistent service supervision and restart. \\
\path{backend/modules/aion_dashboard/actions.py} & Native dashboard teach/ask actions. \\
\path{backend/modules/local_node/local_node_runtime.py} & Native-first AION desktop question route. \\
\path{backend/modules/hexcore/prediction_market_research_engine.py} & Shared verified SD Qwen
runtime-profile resolution. \\
\path{backend/modules/aion_language/query_resonance_api.py} & Corrected legacy resonance query path. \\
\path{backend/modules/aion_language/thesaurus_linker.py} & Deterministic legacy relation export. \\
\path{backend/modules/aion_cognition/advanced_cognition/lexical_task_runner.py} & Compatibility
runner redirected to native examination. \\
\path{backend/tests/test_native_lexical_intelligence.py} & Native English and teacher-separation tests. \\
\path{backend/tests/test_native_english_conversation.py} & Native answer, abstention, evidence,
reference and self-improvement tests. \\
\path{backend/tests/test_native_english_progress.py} & Progress and health projection tests. \\
\path{data/aion/native_english_teacher/} & Small Mac-resident teacher state and audit ledger. \\
\path{/Volumes/TESSARISQ38/TessarisLearning/AION/} & SD-resident lexical database and corpus. \\
\path{/Volumes/TESSARISQ38/TessarisModelPackages/} & SD-resident verified model packages. \\
\bottomrule
\end{longtable}

\section{Conclusion}

AION now has a persistent native English substrate, a fifteen-layer governed curriculum,
sealed free-response advancement tests and a working symbolic conversation path. It can parse
supported questions, retrieve verified meanings, plan an answer, realise grammatical English,
resolve a recent reference, verify evidence and abstain without calling a generative model.
Its own failure and conversation ledgers now feed a governed model-free improvement cycle.

This is not yet unrestricted English intelligence. Its verified fact base and grammar breadth
must grow substantially, and independent real-document transfer remains essential. Its most
important property is architectural honesty: curriculum teaching, native memory, examination,
conversation, self-improvement and capability claims are separate and auditable. In native
conversation mode, Qwen is neither AION's voice nor its ghost-writer; the returned English is
constructed from AION's own structured understanding and explicit language rules.

\end{document}
